\section{Uniformly seizable programmes and unrestricted centres}\label{sec:seizable72}
Programme-state comparisons become exact when a single physical
procedure can recover the programme from one use, uniformly over the
whole family. The state/channel discrimination principle is classical
\cite[Definition~36 and Theorem~37]{WBHK72}; see also the channel-box
formulation \cite[Section~III.A]{WW72}. We state its covering consequence
with the centre classes explicit.

Let $\mathcal S$ be the full state space of a fixed finite direct-sum
matrix algebra. Let $\mathcal A(\omega)$ be the instrument obtained
by feeding a state $\omega\in\mathcal S$ into a fixed CPTP processor
along with its input. Let $\mathcal B$ be a fixed one-use tester
which returns a state in $\mathcal S$ from any legal instrument of
the same interface. Assume
\begin{equation}\label{eq:uniformseizing72}
             \mathcal B(\mathcal A(\omega))=\omega
                         \qquad(\omega\in\mathcal S).
\end{equation}
The tester and processor do not depend on the target, the pair of
hypotheses, the tolerance or the covering centre. They may have a
fixed finite quantum memory. In particular a common programme
representation alone is not this hypothesis.

\begin{theorem}[Exact transfer of family covering numbers]\label{thm:seizablecover72}
Under \eqref{eq:uniformseizing72}, for every $N\ge1$,
\begin{equation}\label{eq:seizablemetric72}
 d_N(\mathcal A(\omega),\mathcal A(\tau))
                =\|\omega^{\otimes N}-\tau^{\otimes N}\|_1.
\end{equation}
The map $\mathcal R=\mathcal A\mathcal B$ on legal memoryless
instrument centres is an idempotent one-use superchannel, and
\begin{equation}\label{eq:seizableretraction72}
 d_N(\mathcal A(\omega),\mathcal R(K))
                 \le d_N(\mathcal A(\omega),K).
\end{equation}
For every subset $S\subseteq\mathcal S$ and every radius $\delta$,
its $N$-copy state covering number with arbitrary centres in
$\mathcal S$ equals the $d_N$ covering number of $\mathcal A(S)$
with arbitrary legal memoryless instrument centres.

In particular, for the rank-bounded state family
$\mathcal S_{\boldsymbol r}\subseteq\mathcal S_{m,n}$ with
$v=\sum_y r_y(2n-r_y)-1$, the latter logarithmic covering number is
\begin{equation}\label{eq:seizablebits72}
       (v/2)\log_2N+v\log_2(1/\delta)+O(1)
\end{equation}
for every fixed $0<\delta_*<2$, uniformly over
$N\ge1$, $0<\delta\le\delta_*$. If $\mathcal A$ sends rational
programme states to rational Choi matrices, the same asymptotic law
holds with rational centres and is attained by the inherited rational
factor code followed by $\mathcal A$. The code preserves programme
block ranks and zeros. No Choi-rank preservation of $\mathcal R$
is asserted.
\end{theorem}
\begin{proof}
Give a common processor $N$ independent copies of the programme,
and use one fresh copy at each device call. All tester operations,
including discarding the unused programmes after a stop, compose into
one common channel. Data processing proves the upper bound in
\eqref{eq:seizablemetric72}. Repeating the one-use seizing procedure
$N$ times and retaining its states gives the reverse bound. This is
the established environment-seizing reduction with its finite-use
normalization made explicit.

To implement $\mathcal R(K)$ on an external input, retain that input
in a quantum memory, use one fresh invocation of $K$ inside
$\mathcal B$, and feed the resulting programme together with the
retained input to $\mathcal A$. This is a fixed one-use pre- and
postprocessing, hence a superchannel. Equation
\eqref{eq:uniformseizing72} implies $\mathcal R^2=\mathcal R$ and
$\mathcal R(\mathcal A(\omega))=\mathcal A(\omega)$.
Replacing each slot of an arbitrary tester against
$\mathcal A(\omega)$ and $\mathcal R(K)$ by this construction
produces a common tester against $\mathcal A(\omega)$ and $K$ with
identical two output states. It uses no more device calls and is
compatible with reference entanglement and stopping. Taking suprema
proves \eqref{eq:seizableretraction72}.

A state cover $\{\tau_i\}$ gives an instrument cover
$\{\mathcal A(\tau_i)\}$ by \eqref{eq:seizablemetric72}.
Conversely, map each arbitrary instrument centre $K_i$ to
$\tau_i=\mathcal B(K_i)$. Equations
\eqref{eq:seizableretraction72} and \eqref{eq:seizablemetric72}
show that every target covered by $K_i$ is covered by $\tau_i$ at
the same radius in the state metric. This proves exact equality of
covering cardinalities, not just a two-hypothesis equivalence.
The centres $\tau_i$ need not obey the rank promises on the targets;
this is why \eqref{eq:uniformseizing72} is imposed on the whole
ambient programme space.

Apply the preparation-state law of
Theorem~\ref{thm:rankentropy68}, the preparation-centre retraction
Theorem~\ref{thm:retraction70}, and its full-error extension
Corollary~\ref{cor:fullrange71}. They give \eqref{eq:seizablebits72}
for arbitrary programme centres. Rationality-preserving $\mathcal A$
maps the rational rank-preserving upper code into legal rational
instrument centres, while the unrestricted lower bound applies to
this smaller class. If $v=0$, zero payload suffices for the known
singleton. The constants can depend on the fixed dimensions, rank
bounds and error cap; they are not uniform as $\delta_*\uparrow2$.
\end{proof}

\begin{corollary}[Pauli-channel descriptions]\label{cor:paulicoding72}
Let $\sigma_0=I,\sigma_1=X,\sigma_2=Y,\sigma_3=Z$ be the Pauli
matrices and let
$\Phi_p(\rho)=\sum_{j=0}^3p_j\sigma_j\rho\sigma_j^*$.
Fix a nonempty public allowed support of cardinality $k$ for $p$,
allowing zero probabilities on that support. With arbitrary legal
qubit-channel centres, its optimal reusable description length is
\[
 \frac{k-1}{2}\log_2N+(k-1)\log_2(1/\delta)+O(1),
 \qquad N\ge1,\quad 0<\delta\le\delta_*<2.
\]
Rational upper descriptions preserve zero probabilities and never
increase the Pauli-channel Choi rank of the target.
\end{corollary}
\begin{proof}
The classical programme $p$ controls which Pauli matrix is applied.
One maximally entangled input and a Bell-basis measurement recover
$p$ exactly. Both operations are fixed over the full probability
simplex, so Theorem~\ref{thm:seizablecover72} applies with $n=1$
and one-dimensional blocks at the allowed positions. The orthogonal
Bell Choi vectors show that the Choi rank equals the number of
positive probabilities. The inherited scalar factor encoder keeps
zeros zero and has rational output. This is a coding consequence of
the standard Pauli/Bell seizing mechanism, not a new discrimination
protocol.
\end{proof}
For contrast, $U=(I+iX+iY+iZ)/2$ is unitary, so its channel has
Choi rank one. Its Bell probabilities are all $1/4$, and
$\mathcal A\mathcal B$ sends it to the uniform Pauli mixture of Choi
rank four. This exact example explains why arbitrary-centre
retraction and rank-preserving target encoding are different claims.
The processor in a programme-state comparison is quantum. Description
length in \eqref{eq:seizablebits72} does not count its Hilbert-space
dimension or give an unknown-channel learning or classical physical
simulation theorem.
